\section{Introducción}
\vspace{-0.15cm}

% contexto, problema, justificación, objetivos

La Organización de las Naciones Unidas para la Educación, la Ciencia y la Cultura (UNESCO) señala que el 40\% de la población mundial no tiene acceso a
educación en una lengua que hable y comprenda con fluidez. Si bien este problema afecta a la población en general, su impacto en las personas cuyas lenguas
maternas corresponden a lenguas de bajos recursos es mucho más drástico  \cite{UNESCO}. Una de las principales dificultades en el ámbito educativo es la persistencia de barreras lingüísticas que limitan el acceso equitativo al aprendizaje a estudiantes cuya lengua materna corresponde a una lengua de bajos recursos lingüísticos, como el quechua o el aimara, debido a la limitada información existente, complejidad morfológica y superposiciones semánticas \cite{Demilie2026,Tafa2025}. 


En torno a esta problemática, la inteligencia artificial (IA) ha comenzado a posicionarse como una alternativa con alto potencial para enfrentar estas limitaciones y se ha venido consolidando como una vía prometedora para mejorar la accesibilidad y personalización del aprendizaje. 
Asimismo, revisiones recientes en el campo de la educación lingüística muestran que las herramientas basadas en IA, incluidos los sistemas
conversacionales y los grandes modelos de lenguaje, han ampliado las posibilidades de apoyo pedagógico en contextos multilingües, aunque todavía existen desafíos significativos en términos de equidad, calidad y estabilidad de los resultados educativos \cite{BeldaMedina2022,Nwafor2026,Nouzri2025,Tebourbi2025}.


En particular, los grandes modelos de lenguaje, \textit{LLMs}, por sus siglas en inglés, tales como ChatGPT o Gemini demuestran un gran desempeño en trabajos que requieren procesamiento del lenguaje natural, sin embargo, su rendimiento en tareas que requieren de un razonamiento más explícito, aún constituye un ámbito inexplorado y con varias inconsistencias. En \cite{Onan2026} se realizó un estudio que comparó seis \textit{LLMs} usados para preguntas de un examen en turco e indonesio, los resultados mostraron que mientras GPT-4 produce las explicaciones de mayor calidad y alineación pedagógica de manera más consistente, Gemini y Qwen 3 logran un desempeño competitivo, pero menos estable en categorías cognitivas, y por su parte DeepSeek-V2 tiene un desempeño sólido en tareas de razonamiento, pero muestra una menor estabilidad temporal, en contraste, Copilot y Phi-2 muestran un menor rendimiento en tareas de pensamiento complejo. Tales hallazgos tienen implicancias directas para la equidad educativa, el rendimiento estudiantil y la educación de calidad en contextos de bajos recursos, especialmente en entornos donde las herramientas basadas en IA son el principal recurso de aprendizaje. 


De manera complementaria, trabajos recientes muestran que estas tecnologías también se están explorando en contextos de aprendizaje de lenguas mediante sistemas multiagente, asistentes conversacionales y flujos de aprendizaje personalizados \cite{Nouzri2025,Tebourbi2025}. Estas evidencias son relevantes para las lenguas de bajos recursos, sugieren que la utilidad educativa de la IA depende tanto de la potencia del modelo como de su adaptación lingüística y pedagógica al contexto de uso.


En este contexto, resulta relevante examinar no solo la capacidad lingüística de los modelos, sino también su utilidad pedagógica y su estabilidad en contextos multilingües. En ese sentido, \cite{Onan2026} señala que una limitación predominante en las investigaciones actuales es el enfoque casi exclusivo en lenguas de alto recurso, tales como el inglés, lo que deja un vacío significativo en la comprensión de cómo estas tecnologías pueden aplicarse en lenguas de bajos recursos. Esto, a su vez, limita el desarrollo de soluciones efectivas y contextualizadas para estos grupos lingüísticos.


De igual manera, en una revisión sistemática orientada a examinar el panorama actual de la investigación sobre IA generativa en la educación del idioma árabe \cite{Moustafa2026}, se evidenció la escasez de recursos académicos luego de recuperar inicialmente 87 registros en bases de datos globales como Scopus y Web of Science, de los cuales únicamente 4 estudios, que representan apenas el 4.6\%, cumplieron con los criterios de inclusión pedagógica y empírica. Aunque se centra en un contexto lingüístico particular, este hallazgo permite inferir que la evidencia disponible sobre el uso de la IA en entornos educativos multilingües todavía es escasa.


A pesar de que la inteligencia artificial en la educación constituye un campo en expansión, la literatura científica disponible todavía se encuentra dispersa. Mientras algunas investigaciones se centran en el desempeño técnico de modelos de traducción, otras analizan experiencias de uso educativo de la IA sin centrarse específicamente en lenguas de bajos recursos ni en la reducción de barreras lingüísticas. Debido a esto, aún no se observan con suficiente claridad los enfoques que han sido más utilizados, qué resultados educativos se han reportado y cuáles son los principales retos metodológicos, éticos y tecnológicos en este campo. 


Dada la dispersión temática y metodológica de la evidencia disponible, resulta necesaria una revisión sistemática de la literatura que permita esclarecer el estado actual del conocimiento, identificar tendencias, oportunidades y vacíos de investigación, así como reconocer las estrategias y metodologías más eficaces, particularmente en contextos educativos donde las tecnologías emergentes pueden contribuir a promover la equidad y la innovación, como también se señala en \cite{Yoshida2025}.


En respuesta a esta necesidad, se plantea una revisión sistemática de la literatura científica sobre el uso de la inteligencia artificial para mejorar la accesibilidad educativa en lenguas de bajos recursos, con especial atención a la reducción de barreras lingüísticas. Esta revisión tiene como propósito sintetizar los enfoques metodológicos empleados, los resultados educativos reportados y los principales desafíos y oportunidades identificados en la literatura, a fin de ofrecer una visión integral del estado actual del campo y orientar futuras investigaciones.
